\begin{table}[t]
\begin{adjustwidth}{0cm}{0cm}
\centering
\resizebox{\ifdim\width>\textwidth\textwidth\else\width\fi}{!}{
\begin{tabular}{lcccccccccccccccc}
\toprule[1.2pt]
                        & \multicolumn{2}{c}{Aspirin} & \multicolumn{2}{c}{Benzene} & \multicolumn{2}{c}{Ethanol} & \multicolumn{2}{c}{Malonaldehyde} & \multicolumn{2}{c}{Naphthalene} & \multicolumn{2}{c}{Salicylic acid} & \multicolumn{2}{c}{Toluene} & \multicolumn{2}{c}{Uracil} \\
\cmidrule[0.6pt]{2-17}
Methods                                               & energy       & forces       & energy       & forces       & energy       & forces       & energy          & forces          & energy         & forces         & energy           & forces          & energy       & forces       & energy       & forces      \\
\midrule[1.2pt]
SchNet~\citep{schnet}                          & 16.0         & 58.5         & 3.5          & 13.4         & 3.5          & 16.9         & 5.6             & 28.6            & 6.9            & 25.2           & 8.7              & 36.9            & 5.2          & 24.7         & 6.1          & 24.3        \\
DimeNet~\citep{dimenet}                          & 8.8          & 21.6         & 3.4          & 8.1          & 2.8          & 10.0         & 4.5             & 16.6            & 5.3            & 9.3            & 5.8              & 16.2            & 4.4          & 9.4          & 5.0          & 13.1        \\
PaiNN~\citep{painn}                                                 & 6.9          & 14.7         & -            & -            & 2.7          & 9.7          & 3.9             & 13.8            & 5.0            & 3.3            & 4.9              & 8.5             & 4.1          & 4.1          & 4.5          & 6.0         \\
TorchMD-NET~\citep{torchmd_net}                                           & \textbf{5.3}          & 11.0         & 2.5          & 8.5          & 2.3          & 4.7          & 3.3             & 7.3             & \textbf{3.7}            & 2.6            & \textbf{4.0}              & 5.6             & \textbf{3.2}          & 2.9          & \textbf{4.1}          & 4.1         \\
NequIP ($L_{max} = 3$)~\citep{nequip}                              & 5.7          & 8.0          & -            & -            & \textbf{2.2}          & 3.1          & 3.3             & 5.6             & 4.9            & \textbf{1.7}            & 4.6              & \textbf{3.9}             & 4.0          & \textbf{2.0}          & 4.5          & \textbf{3.3}         \\

\midrule[0.6pt]

Equiformer ($L_{max} = 2$)                                           & \textbf{5.3}          & 7.2          & \textbf{2.2}          & \textbf{6.6}          & \textbf{2.2}          & 3.1          & 3.3             & 5.8             & \textbf{3.7}            & 2.1           & 4.5              & 4.1             & 3.8          & 2.1          & 4.3          & \textbf{3.3}        \\
Equiformer ($L_{max} = 3$)                                           & \textbf{5.3}          & \textbf{6.6}          & 2.5          & 8.1          & \textbf{2.2}          & \textbf{2.9}          & \textbf{3.2}             & \textbf{5.4}             & 4.4            & 2.0           & 4.3              & \textbf{3.9}            & 3.7          & 2.1         & 4.3          & 3.4        \\
\bottomrule[1.2pt]
\end{tabular}
}
\end{adjustwidth}
\vspace{-3mm}
\caption{\textbf{MAE results on MD17 testing set.} Energy and force are in units of meV and meV/$\angstrom$.}
\vspace{-5.5mm}
\label{tab:md17_result}
\end{table}

\vspace{-3mm}
\subsection{MD17}
\vspace{-3mm}
\label{subsec:md17}


\paragraph{Dataset.}
The MD17 dataset~\citep{md17_1, md17_2, md17_3} (CC BY-NC) consists of molecular dynamics simulations of small organic molecules, and the goal is to predict their energy and forces.
We use $950$ and $50$ different configurations for training and validation sets and the rest for the testing set.
Forces are derived as the negative gradient of energy with respect to atomic positions. 
We minimize MAE between prediction and normalized ground truth.

\vspace{-3.25mm}
\paragraph{Training Details.}
Please refer to Sec.~\ref{appendix:subsec:md17_training_details} in appendix for details on architecture, hyper-parameters and training time.


\vspace{-3.25mm}
\paragraph{Results.}
We train Equiformer with $L_{max} = 2 \text{ and } 3$ and summarize the results in Table~\ref{tab:md17_result}.
Equiformer achieves overall better results across 8 molecules compared to each individual model.
%%%\revision{Compared to TorchMD-NET, which is also an equivariant Transformer, Equiformer obtains better MAE results.} 
%since the proposed equivariant graph attention is more expressive and can support vectors of higher degree $L$ (i.e., we use $L_{max} = 2$) instead of restricting to type-$0$ and type-$1$ vectors (i.e., $L_{max} = 1$).}
%%%\revision{The difference lies in the proposed equivariant graph attention, which is more expressive and can support vectors of higher degree $L$ (i.e., we use $L_{max} = 2 \text{ and } 3$) instead of restricting to type-$0$ and type-$1$ vectors (i.e., $L_{max} = 1$).}
Compared to TorchMD-NET, which is also an equivariant Transformer, the difference lies in the proposed equivariant graph attention, which is more expressive and can support vectors of higher degree $L$ (i.e., we use $L_{max} = 2 \text{ and } 3$) instead of restricting to type-$0$ and type-$1$ vectors (i.e., $L_{max} = 1$).
For the last three molecules, although Equiformer with $L_{max} = 2$ achieves lower force MAE but higher energy MAE, we can adjust the weights of energy loss and force loss so that Equiformer achieves lower MAE for both energy and forces as shown in Table~\ref{appendix:tab:md17_torchmdnet} in appendix.
%%%\revision{Compared to NequIP, which also uses irreps features and $L_{max} = 3$, Equiformer with $L_{max} = 2$ achieves overall lower MAE although including higher $L_{max}$ can improve performance and Equiformer uses smaller $L_{max}$.}
Compared to NequIP, which also uses irreps features and $L_{max} = 3$, Equiformer with $L_{max} = 2$ achieves overall lower MAE although including higher $L_{max}$ can improve performance.
When using $L_{max} = 3$, Equiformer achieves lower MAE results for most molecules.
This suggests that the proposed attention can improve upon linear messages even when the size of training sets becomes small.
We compare the training time of NequIP and Equiformer in Sec.~\ref{appendix:subsec:md17_training_time_comparison}.
Additionally, for Equiformer, increasing $L_{max}$ from $2$ to $3$ improves MAE for most molecules except benzene, which results from overfitting.

%\todo{[Training time of NequIP]}
%\todo{[Equiformer with $L_{max} = 3$]}